% =============================================================================
% Relatório Acadêmico — Geólogo Digital / OpenCode Ecosystem Core
% Orquestrador: /marceloclaro
% Formato: ABNT (NBR 14724, NBR 6023, NBR 10520) — Qualis A1
% =============================================================================
\documentclass[
  a4paper, 12pt, openany, oneside,
  chapterstyle=ABNT, section=TOCentries
]{memoir}

% ── Codificação e Idioma ──────────────────────────────────────────────────
\usepackage[utf8]{inputenc}
\usepackage[T1]{fontenc}
\usepackage[brazil]{babel}

% ── Geometria ABNT (3cm superior/esquerda, 2cm inferior/direita) ──────────
\usepackage{geometry}
\geometry{
  top=3cm, bottom=2cm, left=3cm, right=2cm,
  headheight=1.5cm, headsep=1cm, footskip=1.5cm
}

% ── Fontes ────────────────────────────────────────────────────────────────
\usepackage{mathptmx}        % Times New Roman (ABNT recomenda)
\usepackage[scaled=.92]{helvet}
\usepackage{courier}

% ── Parágrafos e Espaçamento ABNT ─────────────────────────────────────────
\setlength{\parindent}{2cm}
\setlength{\parskip}{0pt}
\SingleSpacing
\OnehalfSpacing*

% ── Numeração Progressiva (ABNT NBR 6024) ─────────────────────────────────
\maxsecnumdepth{subsection}
\settocdepth{subsection}
\renewcommand{\thesection}{\arabic{section}}
\renewcommand{\thesubsection}{\thesection.\arabic{subsection}}

% ── Pacotes Essenciais ────────────────────────────────────────────────────
\usepackage{graphicx}
\usepackage[hyphens]{url}
\usepackage[
  colorlinks=true,
  linkcolor=black,
  citecolor=black,
  urlcolor=blue,
  bookmarks=true,
  hypertexnames=false
]{hyperref}
\usepackage{booktabs}
\usepackage{longtable}
\usepackage{array}
\usepackage{xcolor}
\usepackage{lastpage}
\usepackage{listings}
\usepackage{textcomp}
\usepackage{microtype}
\usepackage{tikz}

% ── Cores institucionais ──────────────────────────────────────────────────
\definecolor{geoverde}{HTML}{3fb950}
\definecolor{geocinza}{HTML}{5c6166}

% ── Cabeçalho e Rodapé ────────────────────────────────────────────────────
\makepagestyle{abnt}
\makeevenhead{abnt}{\small\itshape Geólogo Digital}{}{\small\thepage}
\makeoddhead{abnt}{\small\itshape Geólogo Digital}{}{\small\thepage}
\makeevenfoot{abnt}{}{}{\small\thepage}
\makeoddfoot{abnt}{}{}{\small\thepage}
\pagestyle{abnt}

% ── Fich Catalográfica (verso da folha de rosto) ──────────────────────────
\newcommand{\fichacatalografica}[1]{%
  \newpage
  \thispagestyle{empty}
  \vspace*{\fill}
  \begin{center}
  \fboxsep=12pt
  \fbox{\begin{minipage}[c][6cm]{\textwidth}
    \centering
    {\small
    \begin{tabular}{p{0.12\textwidth}p{0.80\textwidth}}
      \multicolumn{2}{l}{FICHA CATALOGRÁFICA} \\[6pt]
      \textbf{Título:} & VAR_TITULO \\
      \textbf{Data:}   & VAR_DATA \\
      \textbf{ID:}     & VAR_PROJETO_ID \\
      \textbf{Modelo:} & opencode/deepseek-v4-flash-free \\
      \textbf{Tema:}   & VAR_PROMPT \\
    \end{tabular}
    }
  \end{minipage}}
  \end{center}
  \vspace*{\fill}
  \newpage
}

% ── Lista de Símbolos (se houver) ────────────────────────────────────────
\newcommand{\listadesimbolos}{%
  \chapter*{Lista de Símbolos}
  \thispagestyle{empty}
  \begin{longtable}{p{4cm}p{10cm}}
    \texttt{OpenCode} & Plataforma de agentes de IA para engenharia de software \\
    \texttt{/marceloclaro} & Orquestrador primário do ecossistema \\
    SSE & Server-Sent Events — streaming de dados \\
    ABNT & Associação Brasileira de Normas Técnicas \\
    Qualis A1 & Classificação máxima de periódicos CAPES \\
    DOI & Digital Object Identifier \\
  \end{longtable}
  \newpage
}

% ── Resumo ────────────────────────────────────────────────────────────────
\newcommand{\resumoabnt}[1]{%
  \chapter*{Resumo}
  \thispagestyle{empty}
  \OnehalfSpacing
  #1
  \vfill
  \noindent\textbf{Palavras-chave:} geologia, paleontologia, museologia, dissertação acadêmica, OpenCode.
  \newpage
}

% =============================================================================
% INÍCIO DO DOCUMENTO
% =============================================================================
\begin{document}

% ── Pré-Textual ───────────────────────────────────────────────────────────
\frontmatter
\thispagestyle{empty}

% Capa (ABNT NBR 14724)
\begin{tikzpicture}[remember picture, overlay]
  \fill[geoverde!5] (current page.north west) rectangle (current page.south east);
\end{tikzpicture}
\vspace*{3cm}
\begin{center}
{\large MUSEU GEOMAKER --- OPENCODE ECOSYSTEM CORE}\\[0.5cm]
{\large Geólogo Digital --- Orquestrador /marceloclaro}\\[4cm]
{\Huge\bfseries VAR_TITULO}\\[1.5cm]
{\large Dissertação Técnico-Científica}\\[0.3cm]
{\large gerada por inteligência artificial --- OpenCode CLI}\\[4cm]
\vfill
{\large VAR_DATA}\\[0.3cm]
\end{center}
\newpage

% Folha de Rosto
\thispagestyle{empty}
\vspace*{4cm}
\begin{center}
{\Large\bfseries VAR_TITULO}\\[2cm]
{\large Dissertação apresentada ao Geólogo Digital ---}\\
{\large Museu Geomaker --- como resposta à consulta:}\\[1cm]
\begin{quote}
\itshape VAR_PROMPT
\end{quote}
\end{center}
\vfill
\begin{center}
{\large VAR_DATA}
\end{center}
\newpage

% Ficha Catalográfica
\fichacatalografica{VAR_TITULO}

% Lista de Símbolos
\listadesimbolos

% Resumo (gerado dinamicamente)
\resumoabnt{Esta dissertação técnico-científica foi elaborada pelo orquestrador /marceloclaro do OpenCode Ecosystem Core em resposta à consulta do usuário. O conteúdo segue rigor acadêmico com referências reais, DOIs auditáveis e estrutura ABNT Qualis A1.}

% Sumário (ABNT NBR 6027)
\pdfbookmark[0]{Sumário}{toc}
\tableofcontents*
\clearpage

% =============================================================================
% TEXTO PRINCIPAL
% =============================================================================
\mainmatter

VAR_RESPOSTA

\bigskip
\hrule
\bigskip

% ── METADADOS ─────────────────────────────────────────────────────────
\section*{Metadados da Consulta}
\addcontentsline{toc}{section}{Metadados da Consulta}

\begin{longtable}{p{5cm}p{9cm}}
\toprule
\textbf{Parâmetro} & \textbf{Valor} \\
\midrule
\endhead
Modelo & \texttt{opencode/deepseek-v4-flash-free} \\
Orquestrador & \texttt{/marceloclaro} \\
Tokens de entrada & VAR_TOKENS_INPUT \\
Tokens de saída & VAR_TOKENS_OUTPUT \\
Tempo de execução & VAR_ELAPSED~s \\
Custo estimado & R\$~VAR_CUSTO \\
Data da consulta & VAR_DATA \\
Projeto & \texttt{VAR_PROJETO_ID} \\
\bottomrule
\end{longtable}

\bigskip
\noindent{\small\itshape Nota metodológica: as referências bibliográficas e o apêndice de
fichamentos/resenhas críticas apresentados nas seções anteriores foram citados pelo
orquestrador \texttt{/marceloclaro} diretamente no corpo da dissertação, com DOIs e URLs
convertidos em hyperlinks clicáveis para auditoria e verificação (NBR 6023).}

\end{document}
